\documentclass[12pt]{article}
\usepackage[T1]{fontenc}
\usepackage[utf8]{inputenc}
\usepackage{lmodern}
\usepackage{textcomp}
\usepackage{enumitem}
\usepackage{geometry}
\usepackage{graphicx}
\usepackage{amsmath, amssymb}
\geometry{margin=1in}
\begin{document}
\begin{center}
Biology - Undergraduate Level Assessment 1 \
Total Marks: 30 \
Answer all questions.
\end{center}

\section*{Multiple Choice (26 marks)}
\begin{enumerate}[label=\arabic*.]
\item Which of the following facts was established prior to 1859, the year in which Charles Darwin published On the Origin of Species?
\begin{enumerate}[label=(\alph*)]
    \item DNA provides the macromolecular basis of heredity.
    \item Mendelian principles explain why some traits are dominant and others are recessive.
    \item Prokaryotes include two major domains, the Bacteria and the Archaea.
    \item There exist fossilized remains of species that have become extinct.
\end{enumerate}
\item Which of the following would increase the rate at which a gas diffuses between the alveoli of the lung and the blood within a pulmonary capillary?
\begin{enumerate}[label=(\alph*)]
    \item Decreasing the partial pressure gradient of the gas
    \item Decreasing the solubility of the gas in water
    \item Increasing the total surface area available for diffusion
    \item Decreasing the rate of blood flow through the pulmonary capillary
\end{enumerate}
\item Short-term changes in plant growth rate mediated by the plant hormone auxin are hypothesized to result from
\begin{enumerate}[label=(\alph*)]
    \item loss of turgor pressure in the affected cells
    \item increased extensibility of the walls of affected cells
    \item suppression of metabolic activity in affected cells
    \item cytoskeletal rearrangements in the affected cells
\end{enumerate}
\item Targeting of a newly synthesized protein is most likely to require two different signal peptides for which of the following destinations?
\begin{enumerate}[label=(\alph*)]
    \item Plasma membrane
    \item Lysosome
    \item Cytosol
    \item Chloroplast
\end{enumerate}
\item Charles Darwin's proposed conditions for natural selection encompass all of the following with regard to a given population EXCEPT
\begin{enumerate}[label=(\alph*)]
    \item inheritance of both "fit" and "unfit" genes
    \item differential survival and reproductive success
    \item competition for limited resources
    \item overproduction of offspring
\end{enumerate}
\item The first point of entry of water at the roots of a monocot plant is through the cytoplasm of cells of the
\begin{enumerate}[label=(\alph*)]
    \item root cap
    \item Casparian strip
    \item pericycle
    \item endoderm
\end{enumerate}
\item All of the following might be found in connective tissue EXCEPT
\begin{enumerate}[label=(\alph*)]
    \item thrombin
    \item glycosaminoglycans
    \item collagens
    \item fibroblasts
\end{enumerate}
\item Which of the following sets of reactions occurs in the stroma of the chloroplast in plant cells?
\begin{enumerate}[label=(\alph*)]
    \item Calvin cycle
    \item Krebs cycle
    \item Fermentation
    \item Decarboxylation
\end{enumerate}
\item Which of the following most accurately describes a retrotransposon?
\begin{enumerate}[label=(\alph*)]
    \item A DNA sequence that can move from one site in the genome to another without replicating
    \item A DNA sequence that can be deleted from the genome without consequence
    \item A DNA sequence that replicates via an RNA intermediate
    \item A DNA sequence that replicates via a protein intermediate
\end{enumerate}
\item Based on the characteristic population curves that result from plotting population growth of a species, the most effective means of controlling the mosquito...
\begin{enumerate}[label=(\alph*)]
    \item maintain the population at a point corresponding to the midpoint of its logistic curve
    \item opt for zero population control once the K value of the curve has been reached
    \item reduce the carrying capacity cif the environment to lower the K value
    \item increase the mortality rate
\end{enumerate}
\item All of the following statements about plant embryogenesis are correct EXCEPT:
\begin{enumerate}[label=(\alph*)]
    \item The suspensor is derived from the basal cell.
    \item Cotyledons are derived from the apical cell.
    \item Shoot apical meristem formation occurs after seed formation.
    \item Precursors of all three plant tissue systems are formed during embryogenesis.
\end{enumerate}
\item The sight organs of crustaceans and insects contain ommatidia, which make up the individual visual units of the
\begin{enumerate}[label=(\alph*)]
    \item eyespot
    \item simple eye
    \item compound eye
    \item binocular eye
\end{enumerate}
\item Annelids and arthropods are similar to each other in that members of both phyla
\begin{enumerate}[label=(\alph*)]
    \item have segmented bodies
    \item have a closed circulatory system
    \item conduct gas exchange by diffusion through a moist membrane
    \item have well-developed sense organs
\end{enumerate}
\end{enumerate}

\section*{True / False (4 marks)}
\textit{Answer True or False and justify briefly.}
\begin{enumerate}[label=\arabic*.]
\item True or False: The correct answer to 'Keystone species are thought to have profound effects on the structure and composition of ecological communities because they' is 'tend to reduce diversity by eliminating food resources for other species'.
\item True or False: The correct answer to 'Which of the following depicts the correct sequence of membranes of the chloroplast, beginning with the innermost membrane and ending...' is 'Thylakoid membrane, inner membrane, outer membrane'.
\end{enumerate}

\vfill
\noindent\textit{End of Paper}
\end{document}